\section{A concrete repair and a failure of odd-coset monotonicity}
\label{sec:odd-coset-repair}

We continue the $12\times24$ terminal example whose all-odd integer
kernel minimum was certified to be $40$, rather than the sign-vector
value $24$. Exact search finds a one-row exchange, substantially more
efficient than rounding the displayed short odd vector. It also shows
that the odd-coset minimum is not a monotone potential for augmentation:
an explicitly larger frame has no all-odd integer kernel vector at all.
The broader distinction between arithmetic and sign completion is
classical; see Verheiden, \emph{Integral and rational completions of
Hadamard matrices},
\url{https://doi.org/10.1016/0097-3165(80)90028-X}.
The finite certificates and deductions below are independently checked.

Use zero-based coordinates and encode a sign row by the integer having
bit $j$ equal to one exactly when coordinate $j$ is $-1$. Let $P$
have rows with masks
\[
\begin{gathered}
0,16773120,16519104,9656012,13838682,14783914,\\
1944718,11409518,15230494,5921970,3487346,13018356.
\end{gathered}
\]
The rows are orthogonal. The previously studied odd kernel vector is
\[
\begin{split}
v={}&(-1,-1,3,-1,1,1,-3,-1,1,1,-1,1,\\
&1,-1,-1,-1,1,-1,1,-1,1,-1,1,1),
\end{split}
\]
with $Pv=0$ and $\|v\|^2=40$. The two exceptional coordinates are
$2,6$. Its coordinatewise sign rounding is orthogonal to only seven
of the twelve original rows; it violates precisely rows
$\{2,4,6,8,10\}$. Keeping this particular rounded candidate would
therefore require discarding five original rows.

\begin{proposition}[Exact one-row augmentation]
The minimum number of original rows that must be deleted from $P$
to permit an augmenting exchange is exactly one. Specifically delete
row index $1$, with mask $16773120$, and insert rows with masks
\[
3198332,\qquad6782212.
\]
The resulting thirteen rows are orthogonal. More strongly, after this
same deletion the maximum number of mutually orthogonal inserted
rows is exactly five, attained by
\[
3198332,6782212,8735670,11139280,13957672.
\]
The resulting sixteen-row matrix $Q$ is terminal.
\end{proposition}
\begin{proof}
There is no sign row orthogonal to all twelve original rows, so zero
deletions cannot augment. This is certified by exact rational kernel
elimination and enumeration of its $2^{11}=2048$ information-coordinate
sign assignments, up to negation.
Deleting row $1$ gives exactly sixteen antipodal sign-candidate classes.
The checker lists all their masks. Two masks are adjacent exactly when
their Hamming distance is $12$. The five displayed masks form a clique,
and exhaustive recursion on this sixteen-vertex graph finds no
six-clique. Thus five is the exact maximum. All row norms and pairwise
inner products are checked as integers.
If another sign row extended $Q$, it would belong to the same candidate
pool and enlarge the five-clique, which is impossible.
\end{proof}

Deleting row $0$, in contrast, gives only two candidate classes and
no orthogonal pair. Both deletion choices close the odd-coset norm gap
to $24$, since the deleted original sign row itself enters the kernel.
Therefore closing that gap cannot distinguish successful augmentation
from a deletion permitting only replacement.

\begin{proposition}[Augmentation can destroy odd integer feasibility]
The larger frame $Q$ has no all-odd integer kernel vector. In particular
the minimum squared norm in the odd integer coset increases from
$40$ for $P$ to $+\infty$ for $Q$, under an explicit row-augmenting
exchange.
\end{proposition}
\begin{proof}
Let $f=e_0+e_9-e_{11}$. Exact calculation gives
\[
 24\|f\|^2-\sum_{q\text{ row of }Q}(q\cdot f)^2=0.
\]
For orthogonal rows of norm $\sqrt{24}$ this means that $f$ lies in
their rational row span. Its coordinate sum is odd. Any all-odd integer
vector $z$ has $f\cdot z$ odd, whereas $Qz=0$ would force
$f\cdot z=0$. This contradiction proves the assertion and also gives
a small direct certificate of terminality for $Q$.
\end{proof}

\subsection{What coordinatewise rounding does guarantee}
\begin{proposition}[Rounding defect]
Let $P$ have $r$ orthogonal sign rows of length $N$, and suppose
$v\in\ker P$ has entries in $\{\pm1,\pm3\}$. Let $w$ be its
coordinatewise sign, and let $t$ count the coordinates of magnitude
$3$. Define $e_j=w_j$ on those coordinates and $e_j=0$ otherwise.
Then $v=w+2e$ and $Pw=-2Pe$.
If $t$ is odd, $w$ is orthogonal to \emph{none} of the original rows.
If $t$ is even, the number $b$ of rows not orthogonal to $w$ satisfies
\[
 b\le\min\{r,Nt/4\}.
\]
Deleting exactly those $b$ rows permits inserting $w$ as a single
sign row. This need not be an augmenting exchange.
\end{proposition}
\begin{proof}
For every sign row $p$, $p\cdot e\equiv t\pmod2$. If $t$ is odd,
that inner product cannot vanish. If $t$ is even, each nonzero value
has magnitude at least $2$, while Bessel gives
$\sum_p(p\cdot e)^2\le N\|e\|^2=Nt$. Hence $4b\le Nt$.
The identity $Pw=-2Pe$ supplies all remaining claims.
\end{proof}

Thus a short odd vector supplies a deterministic candidate and a
computable deletion set, but its number of exceptional coordinates
alone has not here yielded a bounded augmenting-exchange theorem.
The explicit example has rounding deletion count five and actual
augmentation deletion count one. The earlier sixteen-row order-$24$
example has a certified three-row augmentation barrier and already
fails odd integer feasibility. These examples separate three issues:
parity feasibility, shortness of an odd kernel vector, and mutual
compatibility of enough new sign rows.

Finally, a global column sign change cannot repair any of these
obstructions. Multiplication by a diagonal sign matrix gives a
norm-preserving bijection on sign extensions and on odd integer kernel
vectors, and preserves every row inner product. Column permutations
have the same property. Such operations can simplify a search's
normalization, but cannot change its extension or exchange costs.
